\section{Discussion}
\label{sec:discussion}

\subsection{Key Findings and Their Interpretation}

\textbf{RLHF as a pervasive alignment signal.} Our most striking finding is that RLHF fine-tuning simultaneously shapes 5 of 6 trustworthiness dimensions through a shared optimization signal. The within-family LLaMA-2 natural experiment provides the cleanest evidence: across 7B, 13B, and 70B scale points, RLHF consistently produces joint positive deltas for both safety and ethics, with no counterexamples.

The implication for evaluation practice is immediate: a practitioner who measures 6 dimensions independently is unknowingly measuring the RLHF signal 5 times and the robustness dimension once. The correlation structure suggests that measuring \{truthfulness, fairness, privacy\} provides approximately equivalent ranking power for the RLHF-shaped cluster.

\textbf{Privacy is RLHF-sensitive.} Our original hypothesis classified privacy as RLHF-insensitive alongside robustness. The data strongly refutes this: $\rho(\text{safety, privacy}) = 0.971$ is the strongest pair in the entire 15-pair matrix. Two complementary explanations are plausible: (1) RLHF reward models explicitly penalize privacy-violating outputs alongside harmful outputs; (2) TrustLLM's privacy dimension captures output-level refusal to produce PII, which mechanistically overlaps with safety refusal behavior. Both may be simultaneously true.

\textbf{Robustness isolation as an actionable finding.} Adversarial robustness is the only trustworthiness dimension not co-optimized by RLHF. Robustness improvement requires separate dedicated effort (adversarial training, data augmentation) rather than being a byproduct of standard RLHF alignment.

\subsection{Limitations}

\textbf{Sample size ($n=16$).} With 2 covariates, partial correlation tests have effective $df=12$ ($n - k_\text{covariates} - 2 = 16 - 2 - 2$; see Algorithm~\ref{alg:partial_spearman} note in Section~\ref{sec:method}), limiting power to detect moderate correlations ($|\rho| < 0.55$). The safety-robustness null finding ($\rho=-0.188$, $p=0.519$) is almost certainly a power issue: scale-only ablation ($\rho=-0.771$, $p=0.0008$) confirms the effect exists. Replication with $n \geq 30$ is needed.

\textbf{RLHF mechanism is proposed, not causally proven.} The RLHF explanation is the most parsimonious given within-family evidence, but we cannot rule out alternative confounders in observational cross-model data.

\textbf{TrustLLM-specific operationalization.} Findings are specific to TrustLLM's 2024 evaluation framework. HELM cross-framework replication is identified as important future work.

\textbf{Machine\_ethics MST edge ambiguity.} The machine\_ethics--privacy edge has 68.8\% bootstrap frequency at $n=16$. The minimum evaluation set \emph{size} (3 dimensions) is stable; machine\_ethics topology is ambiguous.

\subsection{Broader Impact}

\textbf{For benchmark design.} Future trustworthiness benchmarks should report a correlation structure summary alongside dimension scores, enabling practitioners to identify redundant dimensions within each model generation.

\textbf{For deployment evaluation.} Our MST-derived minimum evaluation set \{truthfulness, fairness, privacy\} enables principled evaluation compression: 3 dimensions instead of 6 for the correlated cluster, with adversarial robustness measured independently.

\textbf{For alignment research.} The privacy-safety coupling ($\rho=0.971$) suggests RLHF alignment may achieve broader trustworthiness benefits than explicitly targeted, with implications for how constitutional AI and RLHF policies are designed.

\textbf{Potential negative uses.} Understanding which dimensions are correlated could theoretically help adversarial actors target models that score high on correlated dimensions while performing poorly on robustness. Making the structure explicit enables defensive countermeasures (targeted robustness evaluation) that are not possible without this knowledge.
